% ============================================================
%  Chapter 08 — Institutions & Governance · Law & Conscience ·
%               Corporate Governance · Ethics in International
%               Relations · Just War Theory
% ============================================================
\chapter{Institutions, Law, Conscience, Corporate Governance \& Ethics in IR}
\label{ch:institutions-law-conscience-corporate-ir}

\section{Ethical Concerns \& Dilemmas in Public and Private Institutions}

\subsection{Foundational Backdrop --- The Social Contract}
Before discussing institutions, recall the idea that binds the citizen to the state and limits its power.

\begin{KeyConcept}
\begin{itemize}
\item Social contract: citizens implicitly agree to submit to a powerful authority (the state) which, in return, protects them and secures a peaceful society.
\item Constitutionalism: this authority is deliberately kept a limited authority --- the state is powerful but bound, so it cannot use its power for personal ends. Limited authority is the essence of constitutionalism.
\end{itemize}
\end{KeyConcept}

\subsection{Public Administration vs Private Administration}
The core distinction lies in their ultimate objective. Everything else --- accountability, decision-making, dilemmas --- flows from this.

\begin{KeyConcept}
\begin{itemize}
\item Public administration ultimately exists to serve the public interest. Whether it earns a profit or not, it must serve the public interest.
\item Private institutions exist for profit-making. Profit is a central parameter in their decision-making. Making a private company profitable is itself an ethical obligation --- the employee is paid a salary out of that profit.
\end{itemize}
\end{KeyConcept}

\paragraph{UPSC Case Study --- Manager of a Hospital during COVID-19.} A UPSC case study in two parts, first as a public hospital, then as a private hospital.

\begin{PYQLink}
\begin{itemize}
\item Setting: You are the manager of a hospital during COVID-19. COVID-19 is a highly infectious disease; resources are limited and human resources (doctors, nurses, staff) are limited and unlimited doctors do not exist. Staff can catch the infection and even lose their lives; once a COVID case appears, that part of the apparatus must be sealed, causing loss to the whole hospital.
\item Part 1 (public hospital): Knowing fully well that resources are limited and the disease can threaten your staff, what criteria and justification would you give for treating COVID-19 patients --- why put your people in harm's way?
\item Part 2 (private hospital): If you were the manager of a private hospital instead, how would your criteria and justification change?
\end{itemize}
\end{PYQLink}

\begin{ExampleBox}
\textbf{Answer direction}
\begin{itemize}
\item Public hospital: the overriding justification is that you must uphold the public interest --- a public institution cannot refuse the sick.
\item Private hospital: merely treating patients is not sufficient. Survival of the institution also matters, because keeping a private hospital profitable is an ethical obligation and you must strive for profit.
\item A government can print more money, borrow more money and has sovereignty; a private institution's survival rests on its own resources. Hence the reasoning legitimately differs between the two.
\end{itemize}
\end{ExampleBox}

\paragraph{UPSC Case Study --- Faster Vaccine Development in a Private Firm.}
\begin{PYQLink}
\begin{itemize}
\item Setting: You work in a private drug/vaccine firm. Developing the vaccine faster gives a first-mover advantage; the matter is to be reported to the public authority.
\item Key ethical principle: Ethics does not mean inaction, and it does not mean placing obstacles (``rode atkana'') in the path of legitimate work.
\end{itemize}
\end{PYQLink}

\begin{CriticalPoint}
\begin{itemize}
\item A good answer must actively suggest how the vaccine can be developed faster and correctly --- not merely warn against wrongdoing.
\item You should reject wrong research and negative consequences, but you must also propose how the right work can be done, because contributing to the organisation is your ethical obligation. Inaction lets people suffer.
\item Ethics = solution-oriented + effective. The ideal answer shows you are ethical, effective, a problem-solver, and loyal to the private company --- all while remaining ethical and legal.
\end{itemize}
\end{CriticalPoint}

\paragraph{Key Differences --- Public vs Private.}

\begin{tabularx}{\textwidth}{@{}>{\bfseries\color{Accent}}p{2.6cm} Y Y@{}}
\toprule
\rowcolor{AccentDark}\thead{Dimension} & \thead{Public Administration} & \thead{Private Administration} \\
\midrule
Objective & Serve the public interest (wide, diverse: farmers, drivers, etc.) & Give the consumer what they want; make profit (narrow purpose) \\
Accountability & Elaborate --- citizens, media, NGOs, judiciary; sense of accountability much higher & Narrow set --- consumers, consumer rights, shareholders, investors, employees \\
Ability to differentiate & Cannot differentiate between those who can pay and cannot pay; liable to all & Serves consumers; can differentiate; not really liable to those who cannot pay \\
Political context & Operates in a political context; must weigh politically sensitive decisions & Largely free of political interests; decisions taken on own criteria \\
Decision-making speed & Slower --- many rules \& regulations to follow & Faster --- comparatively fewer rules apply; more autonomy \\
Salary / autonomy & Same salary for all (e.g.\ IAS); government cannot make distinctions & Decides salaries; can differentiate; has latitude/autonomy \\
Ethical dilemmas & More --- dealing with multiple, competing interests & Fewer --- relatively fewer stakeholders and interests to pursue \\
\bottomrule
\end{tabularx}

\begin{ExampleBox}
\begin{itemize}
\item Bringing private stars into public administration: big names from the private sector, used to quick decision-making, often fail in public administration, where every decision must be checked for political sensitivity.
\item Laptop procurement: a student wanted to procure laptops for a government office during an Amazon Diwali sale (cheaper), but could not --- a tender and full process were mandatory. Spending public money triggers heavy accountability mechanisms, whereas a private institution would simply buy from Amazon at once.
\end{itemize}
\end{ExampleBox}

\paragraph{Applied Question --- Changing Nature of Institutions \& of Dilemmas.}
\begin{PYQLink}
Applied (not yet asked, anticipated): How will the changing nature of public and private institutions change the nature of ethical dilemmas and ethical concerns?
\end{PYQLink}

\textbf{Nature of dilemma}
\begin{itemize}
\item Private institution: profit is always a central angle --- e.g.\ heavy environmental-regulation compliance lowers profit, so the firm must balance sustainability of profit against workers' safety and wages. The dilemma is tied largely to profit.
\item Public institution: the dilemma is far wider, because public interest itself is wide --- tribal welfare, development, equity, efficiency, accountability, justice, DPSP, fundamental rights --- obligations that constantly come into conflict.
\end{itemize}

\textbf{Nature of ethical concern}
\begin{itemize}
\item Public administration $\rightarrow$ red-tapism. Because of heavy rules and regulations, decision-making is slow, raising the chance of red-tapism (\emph{laal feeta shaahi}, ``red tape''). Max Weber's stress on written rules and regulations slows decisions further.
\item Private administration $\rightarrow$ discrimination. Because of high autonomy, the concern shifts to discrimination --- e.g.\ the glass ceiling: an employer can pay a woman less and a man more and simply claim a ``merit'' difference. High autonomy enables gender discrimination.
\item In public administration the ``glass ceiling'' is generally not heard of, because pay is governed by law --- a woman IAS and a man IAS are paid equally.
\end{itemize}

\subsection{Changing Values under Globalisation}
The neat split (public interest = public bodies; profit = private bodies) is blurring in both directions.

\begin{GovtPolicy}
\begin{itemize}
\item Private firms adopting public interest: CSR is now mandatory for eligible companies, so private companies are adopting public-interest goals.
\item Public bodies adopting profit: PSUs are classified as Maharatna, Miniratna and Navratna. Each PSU has profit and revenue targets; depending on target-fulfilment they are granted more autonomy in decision-making. Government even considers winding up loss-making PSUs --- so profit is becoming important for public bodies too.
\end{itemize}
\end{GovtPolicy}

\subsection{Note --- Conflict of Interest vs Ethical Dilemma}
\begin{CriticalPoint}
\begin{itemize}
\item Ethical dilemmas do not necessarily involve a private interest. Example: the government building a dam has no private interest but is genuinely conflicted between environment and development; obligations may be immediate vs long-term.
\item In a lockdown the dilemma is GDP/economy vs people's lives --- again no personal interest.
\item Broader category: which ethical values you pursue (no personal interest). Narrow category: conflict of interest, where a personal interest is involved.
\end{itemize}
\end{CriticalPoint}

\begin{QuickRevision}
Public administration serves public interest; private serves profit --- and this drives every other difference.\\
Accountability, dilemmas and rules are greater in public bodies; autonomy and speed are greater in private bodies.\\
Concerns diverge: red-tapism in public bodies vs discrimination / glass ceiling in private bodies.\\
Globalisation is blurring the split --- CSR pulls private firms toward public interest; Ratna targets pull PSUs toward profit.
\end{QuickRevision}

\subsection{Ethical Imperative in Public Administration}
This asks: why should public administration follow ethical standards? These arguments were built earlier (foundational values for civil services; consequences of ethics for an organisation) and are merely restated here.
\begin{itemize}
\item Trust and confidence --- e.g.\ the Election Commission as an institution that runs on public trust.
\item Effective delivery --- e.g.\ Direct Benefit Transfer (DBT) shows why ethical conduct matters for outcomes.
\item Ethical standards such as integrity, objectivity and non-partisanship are themselves the answer to ``why be ethical.''
\end{itemize}

\begin{KeyConcept}
\begin{itemize}
\item Two-layer argument technique: for any ethical value, combine a generic argument (why any organisation should be ethical) with a specific argument for that value.
\item Non-partisanship example: generic layer --- corruption, efficiency, legitimacy, accountability; specific layer --- three levels of impact: (i) bureaucracy's motivation and morale; (ii) transfers, postings, promotions on merit vs party interest; (iii) the citizen and political parties. Generic + specific = a high-scoring answer.
\end{itemize}
\end{KeyConcept}

\subsection{Ethical Challenges in Public Administration}
These are the problems public administration faces because of the mixed and often excessive expectations placed on it. Whatever was earlier discussed as ``negativity'' or ``concerns'' in government becomes an example of an ethical challenge here.
\begin{itemize}
\item Red-tapism; rigid hierarchy; poor team-building; top-down implementation.
\item Well-meaning neutrality becoming apathy; rules becoming ends in themselves; lack of coordination; prevailing corruption.
\item Violations of non-partisanship; discrimination; lack of integrity; limitations of objectivity; plus the limitations of the Max Weber model.
\end{itemize}

\begin{ExampleBox}
\begin{itemize}
\item West Bengal: bureaucrats appointed as Chief Secretary and Chief Advisor across governments.
\item Bureaucrats resigning from service to immediately join a political party and contest elections --- a violation of non-partisanship.
\end{itemize}
\end{ExampleBox}

\begin{CriticalPoint}
\begin{itemize}
\item The listed ethical challenges to reproduce in an answer: corruption, political interference, lack of accountability, lack of transparency, weak ethical training, cultural factors.
\item The challenges of government work culture are the same set of challenges; a lack of emotional intelligence can also be cited as an example.
\end{itemize}
\end{CriticalPoint}

\begin{QuickRevision}
Ethical imperative = why to be ethical (trust, delivery, integrity/objectivity/non-partisanship) using generic + specific arguments.\\
Ethical challenges = what goes wrong (red-tapism, apathy, corruption, political interference, weak training, cultural factors, non-partisanship violations).
\end{QuickRevision}

\section{Law as a Source of Ethical Guidance}
Recall the earlier ``Sources of Ethics'' diagram: how does a person know what is right and wrong? Multiple sources tell us --- our religion, the law, our culture. This topic examines in detail how law guides us on right and wrong.

\subsection{Importance of Law}

\paragraph{Codification of Public Morality.}
\begin{KeyConcept}
\begin{itemize}
\item Law is the codification of public morality through the legislature. In a democracy: we hold elections $\rightarrow$ elected representatives debate $\rightarrow$ they declare what is legal and illegal.
\item What society collectively believes to be right and wrong is expressed by its representatives and turned into law. The overlap is not total --- not every public morality becomes law --- but law, at least to some extent, expresses what society desires through its representatives.
\end{itemize}
\end{KeyConcept}

\begin{GovtPolicy}
\begin{itemize}
\item Nirbhaya case: public demand led to the Criminal Law Amendment; the definition of sexual harassment was widened and the death penalty was provided for certain acts. Because juveniles were going ``scot-free'' (lenient punishment), the Juvenile Justice Act was amended so that 16--18 year-olds committing heinous crimes can be tried as adults.
\item Online betting: as online betting spread (suicides, unemployment), the government moved to ban it --- new social problems being codified into law.
\end{itemize}
\end{GovtPolicy}

\paragraph{Fighting Social Evil.} Law identifies and fights social evils --- e.g.\ child labour, child marriage, dowry --- ensuring how such evils are to be combated.

\paragraph{Delivering Social Justice.} Through reservation, affirmative action and EWS, law becomes an instrument of social justice.

\paragraph{Adding New Ethical Standards.}
\begin{DataFact}
Because law is enforceable, it can add entirely new ethical standards that society did not previously observe. India became the first country to legally mandate CSR --- companies did not do CSR before; the law added a new ethical standard.
\end{DataFact}

\paragraph{Business of Rights \& Protecting Dignity.} Rights become the ``business'' of law --- Fundamental Rights, MGNREGA, RTI, RTE. Law protects your dignity: no discrimination on the basis of caste, religion or descent; no forced conversion. Instruments: Protection of Civil Rights Act; Manual Scavenging Act.

\paragraph{Promoting Good Governance.} Law promotes good governance --- e.g.\ the Jan Vishwas Act.

\begin{CriticalPoint}
\begin{itemize}
\item Laws and regulations are a source of ethical guidance because they tell you clearly what is right and what is wrong, and they codify public morality through representation.
\item They represent the spirit of society --- what we collectively believe to be right and wrong (never a perfect overlap, but a genuine expression). Examples: Prevention of Corruption Act, Prevention of Atrocities Act.
\end{itemize}
\end{CriticalPoint}

\begin{QuickRevision}
Six roles of law: codify public morality, fight social evil, deliver social justice, add new ethical standards, protect rights \& dignity, promote good governance.
\end{QuickRevision}

\subsection{Limitations of Law}
Law is a source of ethical guidance, but an incomplete one. The central limitation: something can be technically legal but not ethical.

\paragraph{Legal but Not Ethical.}
\begin{ExampleBox}
\begin{itemize}
\item Surrogate advertisement: promoting branded alcohol is not allowed, so firms bypass the law by advertising packaged mineral water under the same brand. The law has been bypassed --- technically legal, but the intent is unethical.
\item Marital rape: it is not clearly illegal in India, so the law fails to give ethical guidance in such cases --- the law does not cover it even though the act is against dignity.
\end{itemize}
\end{ExampleBox}

\paragraph{Enforcement Failures \& Evasion.}
\begin{itemize}
\item Lack of enforcement --- dowry: dowry is prohibited yet widely practised, because enforcement agencies do not enforce the law. ``The strength of society lies in law, but not in the morality of the people.'' Law has failed to be a source of ethical guidance.
\item Evasive tendency --- Base Erosion and Profit Shifting (BEPS): MNCs exploit treaties, showing operations in low-tax countries to minimise tax. They misuse legal provisions and loopholes --- legal but unethical, working against the spirit of the law.
\end{itemize}

\paragraph{Negative Perception --- No Internal Change.}
\begin{CriticalPoint}
\begin{itemize}
\item Law is an external source of guidance, working through punishment and deterrence (``don't be corrupt or you'll go to jail''). It does not necessarily change a person's behaviour or internal character.
\item Where the law cannot reach --- e.g.\ the image of people looting a fallen tanker, or everyday corruption --- a person of poor internal character will still act wrongly. Law does not transform the inner self.
\end{itemize}
\end{CriticalPoint}

\paragraph{Law is Minimum Morality \& Involves Discretion.}
\begin{KeyConcept}
\begin{itemize}
\item Law is minimum morality, not maximum morality --- not every ethical standard is turned into enforceable law. Fundamental Duties are ethical standards but are not all legally enforceable.
\item Law cannot envisage every situation, so it must leave discretion. Example: Article 19 reasonable restrictions --- what counts as ``reasonable'' is left to discretion. Hence law cannot always provide ethical guidance in all situations.
\end{itemize}
\end{KeyConcept}

\begin{PYQLink}
UPSC: ``The rules and regulations provided to the civil servants are the same, yet positive-minded officers achieve much more than negative-minded officers.'' --- Positive-minded officers use the law in its positive spirit: they interpret reasonable restrictions correctly, do not file false sedition cases, do not resort to censorship. Because they are ethical, they use the same discretion for good.
\end{PYQLink}

\begin{QuickRevision}
Limitations: legal $\neq$ ethical (surrogate ads, marital rape); lack of enforcement (dowry); evasion / loopholes (BEPS); no internal change (external deterrence only); law = minimum morality; and unavoidable discretion (Art.\ 19 reasonable restrictions).
\end{QuickRevision}

\section{Conscience as a Source of Ethical Guidance}
Because law leaves so much discretion, conscience becomes crucial in filling the gap.

\subsection{What Conscience Is}
\begin{KeyConcept}
\begin{itemize}
\item Conscience is our cognitive act of mind --- our inner sense of right and wrong.
\item It is a cognitive act (using our cognition), not merely a feeling. It is an inner sense, an internal compass, but ultimately a cognitive judgment in which reasoning is involved.
\item It is shaped by experience --- our life experiences, whatever we have lived and learnt, our teachers, the society we live in, the religion we follow, our parents and family. Everything contributes; conscience is a reflection of all those contributions.
\end{itemize}
\end{KeyConcept}

\subsection{How Conscience Becomes a Source of Guidance}
When external sources do not clearly say what is right --- e.g.\ what ``reasonable restriction'' means --- one's own sense of morality guides on right and wrong.

\begin{ExampleBox}
\begin{itemize}
\item Marital rape: the law does not call it illegal, but conscience forbids it as against dignity.
\item Emergencies / disaster management: law gives no clear guidance, so conscience tells the officer what must be done.
\item UBDT (Unbanked Direct Blood Transfer): conscience says this woman cannot be allowed to die --- even if a rule must be broken to save her.
\end{itemize}
\end{ExampleBox}

\subsection{Conscience Demands Personal Responsibility}
\begin{CriticalPoint}
\begin{itemize}
\item Conscience is an inner moral compass and a reflection of values; ultimately it asks you to take personal responsibility for your judgment. You can never hide behind the law.
\item President's visit / VVIP protocol: during a VVIP movement a traffic protocol was imposed and a woman died; the officer pleaded ``the traffic law was applied, what could I do?'' But conscience should have questioned whether that rule outweighed a life.
\item PDS officer: a woman lacking documents --- conscience must weigh the rule against her survival.
\end{itemize}
\end{CriticalPoint}

\begin{ExampleBox}
Gandhiji \& Chauri Chaura: his conscience insisted that a mass movement must remain non-violent, so he called off the movement. A classic example of conscience guiding decision-making.
\end{ExampleBox}

\subsection{Conscience vs Law}

\begin{tabularx}{\textwidth}{@{}>{\bfseries\color{Accent}}p{2.6cm} Y Y@{}}
\toprule
\rowcolor{AccentDark}\thead{Aspect} & \thead{Conscience} & \thead{Law} \\
\midrule
Location & Inner source of guidance & External source of guidance \\
Control & You have control; shaped by upbringing and changeable through new experience & In others' hands; enforced upon you \\
Mode & Cognitive judgment of right and wrong & Punishment / deterrence \\
\bottomrule
\end{tabularx}

\subsection{Can You Always Rely on Conscience? (Is it an Infallible Guide?)}
\begin{PYQLink}
UPSC: Is conscience an infallible guide to ethical decision-making? Can you always rely on it? --- No.
\end{PYQLink}

\begin{CriticalPoint}
\begin{itemize}
\item Conscience is highly subjective and does not necessarily produce ethical decisions.
\item Honour killing: parents believe they are right, ``saving their honour.'' Kasab / suicide bombers: felt no regret because they were trained that way.
\item Conscience can be unethically trained --- the internal sense of right and wrong can itself be wrong. Experiences and learnings can be narrow; religious doctrine can be absorbed uncritically.
\end{itemize}
\end{CriticalPoint}

\begin{KeyConcept}
\begin{itemize}
\item Therefore conscience should be trained in constitutional morality.
\item Socrates: ``The unexamined life is not worth living'' --- we need habitual self-reflection / introspection to train the conscience better.
\item Abraham Lincoln: ``Nothing in this world is considered wholly good or wholly bad.'' Knowledge needs constant churning and critical thinking; what is just in one context can become unjust in another.
\end{itemize}
\end{KeyConcept}

\begin{ExampleBox}
\begin{itemize}
\item HDFC Bank chairman resigned citing conscience --- he could no longer accept certain practices.
\item Kasab's lawyer initially refused to defend him --- a matter of conscience.
\item Parents paying for paper leaks --- a conscience issue.
\end{itemize}
\end{ExampleBox}

\subsection{UPSC PYQ --- Max Weber on Bureaucratic Morality}
\begin{PYQLink}
UPSC (Max Weber): ``It is not wise to apply to public administration the sort of moral and ethical norms we apply to matters of personal conscience. It is important to realise that state bureaucracy might possess its own independent bureaucratic morality.''
\end{PYQLink}

\paragraph{First half --- Weber's point.}
\begin{itemize}
\item The norms we apply to our personal conscience cannot run a government. A bureaucrat must be impersonal.
\item Example: a Jain who believes in non-violence, if appointed head of a fishery board / food-processing body (meat value-addition, poultry, fisheries), cannot administer if he imports his personal conscience. Personal conscience would also push him to help his brother and friends $\rightarrow$ nepotism. Hence Weber's value-neutrality, objectivity, impersonality.
\end{itemize}

\begin{KeyConcept}
Independent bureaucratic morality = value-neutrality, impersonality, legal-rational authority, hierarchy, division of labour, career as a profession, and adherence to written rules and regulations. If you are part of the bureaucracy, follow these; keep personal matters out of the equation.
\end{KeyConcept}

\paragraph{Second half --- but conscience is still needed.}
\begin{CriticalPoint}
\begin{itemize}
\item If you keep your own morality entirely out, the result is apathy, rules becoming ends, and red-tapism.
\item Weber says follow laws, rules and regulations --- but laws do not tell you what is right and wrong (e.g.\ what ``reasonable restriction'' means). So you must, at times, listen to conscience, especially in emergencies and where lives are at stake.
\item Conclusion: ultimate responsibility is yours; but conscience has limits and can be wrong, so all of us should train our conscience in constitutional morality and practise reflective thinking.
\end{itemize}
\end{CriticalPoint}

\subsection{Crisis of Conscience}
\begin{KeyConcept}
\begin{itemize}
\item Crisis of conscience is the deep sense of guilt and crisis that arises when you firmly believe something is right (your conscience says so) yet act against it. That immense guilt --- even from a mistake made knowingly or unknowingly that you regard as wrong --- is a crisis of conscience.
\item Conscience is an important sanction mechanism (the ``fifth sanction''): acting far from one's conscience causes psychological suffering / mental pain.
\item Whistleblower: if you fail to blow the whistle and many people die because of your inaction, the later guilt and shame is a crisis of conscience --- which is why conscience works as a sanction.
\end{itemize}
\end{KeyConcept}

\begin{ComparisonBox}
Crisis of conscience vs cognitive dissonance: choosing a blue vs a pink phone can create guilt / cognitive dissonance, but not necessarily a crisis of conscience. A crisis of conscience is always about ethics --- the ethical angle is the key difference between cognitive dissonance and a crisis of conscience.
\end{ComparisonBox}

\begin{QuickRevision}
Conscience = inner, cognitive sense of right/wrong, shaped by experience; fills the discretion left by law.\\
It demands personal responsibility --- never hide behind rules.\\
But it is not infallible (can be unethically trained) $\rightarrow$ train it in constitutional morality with habitual self-reflection.\\
Weber PYQ: keep personal conscience out for impersonality, yet conscience is indispensable where law is silent.\\
Crisis of conscience = guilt from acting against one's conscience; distinguished from cognitive dissonance by the ethical angle.
\end{QuickRevision}

\section{Corporate Governance --- Ethics in Private Institutions}

\subsection{Why We Study This Topic (The Core Argument)}
\begin{KeyConcept}
\begin{itemize}
\item The private sector pursues profit and often assumes that being ethical and being profitable are opposed --- controlling pollution, paying fair wages, ensuring fire safety, working honestly all raise cost and seem to cut profit. Believing this, some firms resort to unethical means.
\item The single purpose of this topic is to be able to argue that being ethical actually helps good business --- ``good ethics is good business.'' This exact argument is what you will use in case studies and answer-writing.
\end{itemize}
\end{KeyConcept}

\subsection{What Corporate Governance Is}
\begin{itemize}
\item When you buy a share (Jio, Reliance, Tata, etc.) you become an owner --- which requires trust.
\item Management (the CXOs) runs the company day-to-day using shareholders' funds as a trustee, and therefore has an ethical obligation to use those funds ethically for the company's growth.
\item The Board of Directors supervises the functioning; regular board meetings approve plans.
\end{itemize}

\begin{KeyConcept}
\begin{itemize}
\item Corporate governance is the fairness of company management --- being fair to every stakeholder, not just shareholders.
\item Stakeholder = whosoever is impacted by the operations of the company: owners/shareholders, consumers, employees, suppliers, competitors, government/regulators and society.
\item The firm must navigate between profitability, morality, legality and all social \& environmental obligations --- and we argue that by being environmentally sustainable, legal, moral and socially conscious, a company becomes more profitable in the long run.
\end{itemize}
\end{KeyConcept}

\subsection{Obligations Towards Shareholders}
\begin{itemize}
\item Do not misuse funds for private purposes. Money raised for the company must not be diverted to personal use.
\end{itemize}

\begin{ExampleBox}
Blue Smart Cab: the sole promoter (Jain) raised investors' money and bought apartments/flats worth crores for himself --- a clear violation, since the money came for the company, not private use.
\end{ExampleBox}

\begin{itemize}
\item Ensure a fair and sustainable return (at least the intention to). People invest for growth and returns; having no intention to return their money is a form of deceit.
\item Protect minority shareholders. Just as society has majoritarianism, the shareholding pattern has majoritarianism --- a large holder (e.g.\ Reliance) nominates most of the board and can control $\sim$95\%, so minority-shareholder interest can be compromised.
\end{itemize}

\begin{GovtPolicy}
Companies Act provision: at least one-third of directors must be independent directors --- directors with no connection to company management, who independently verify that the interest of minority shareholders is not compromised.
\end{GovtPolicy}

\begin{itemize}
\item Provide accurate and fair information. When a company brings an IPO it advertises its performance; it must give accurate, fair information so investors make an informed decision. Shareholders receive updates, may participate in the general meeting and may vote.
\end{itemize}

\begin{ExampleBox}
Satyam scam: the company faked data to raise more money --- violating the duty of accurate information so that investors could not decide in an informed way. Ultimately the truth came out and the business collapsed.
\end{ExampleBox}

\begin{DataFact}
\begin{itemize}
\item Good ethics = good business (shareholders): timely repayment builds credibility $\rightarrow$ cheaper loans and better valuations (higher company value); transparency attracts investors.
\item HDFC Bank: on the chairman's resignation the bank lost lakhs of crores in a single day and its valuation dropped --- public perception of unethical conduct causes business loss.
\item IndusInd Bank: an accounting-mismatch news item wiped $\sim$27\% off the share in one day. Yes Bank is a similar example.
\end{itemize}
\end{DataFact}

\subsection{Obligations Towards Consumers}
\begin{itemize}
\item Uphold consumer rights --- including the right to safe products, right to be informed, and other consumer rights.
\end{itemize}

\begin{PYQLink}
UPSC case study: a reputed food company licensed to sell food products in the domestic market was selling unsafe/substandard products. Ethical issue: consumer rights violated --- specifically the right to safe products. The firm failed precisely because it did not value its consumer.
\end{PYQLink}

\begin{KeyConcept}
Valuing the consumer $\rightarrow$ word of mouth $\rightarrow$ brand value, market loyalty, lower marketing spend $\rightarrow$ more sales $\rightarrow$ profit. Ethical service can itself be a source of distinction from competitors.
\end{KeyConcept}

\begin{ExampleBox}
\begin{itemize}
\item Tata Motors --- cars seen as safe $\rightarrow$ brand value $\rightarrow$ lower marketing expense. iPhone --- valued for privacy $\rightarrow$ people buy it.
\item Certifications signalling ethics: ISI Mark (safety); Organic certificate; Kimberley Certification (diamond is ethically sourced, not a ``blood diamond'' financing murder); ethically-sourced cocoa/chocolate; ``no refined oil'' / vegetable-oil claims; Biotique animal-cruelty-free products; ``baked, not fried.''
\item Patagonia donates 1\% of sales to the environment. A CP restaurant (``Indus Flavour'') advertised ``we prepare fresh curries'' --- ethical practice used as a point of distinction.
\end{itemize}
\end{ExampleBox}

\subsection{Obligations Towards Employees}
\begin{itemize}
\item Merit-based appointment, promotion and pay; no favouritism, casteism or gender bias.
\item Will not discriminate on the basis of race, caste, sex, etc.; will not indulge in exploitation (no child labour); ensure work-life balance, overall well-being and mental health.
\item The glass ceiling (women not promoted) is wrong.
\end{itemize}

\begin{ExampleBox}
\begin{itemize}
\item Google: rewards employees handsomely, trusts them, does not micromanage, offers flexibility and even nap pods $\rightarrow$ attracts merit; people compete to work there.
\item Tata: after employee deaths during COVID, Tata supported the families with education and financial assistance --- building deep trust and faith.
\item Negative: the Noida protest; the TCS Nashik/Pune case (sexual harassment with no hearing).
\end{itemize}
\end{ExampleBox}

\begin{KeyConcept}
Fairness to employees reduces attrition (saving recruitment and training cost) and yields better motivation, higher productivity and long-term sustainability --- i.e.\ higher efficiency. Good ethics is good business.
\end{KeyConcept}

\subsection{Obligations Towards Government / Regulators}
Why should a company comply with regulations? Note that ``to avoid penalty'' is a weak ethical argument --- like saying ``don't steal because you'll go to jail.'' The stronger arguments follow.
\begin{itemize}
\item Avoid legal cost --- penalties; non-compliance can shut operations. (A real reason, but not a deep ethical one.)
\item Gain credibility and the trust of consumers --- complying with regulations (e.g.\ bank regulations) allows consumers to trust you.
\end{itemize}

\begin{GovtPolicy}
\begin{itemize}
\item Risk management: profit-making involves risk; regulations help manage it. In a bank, deposits are liabilities and loans are assets; reckless (greedy) lending creates NPAs (Non-Performing Assets).
\item RBI regulations require a buffer mechanism (Tier-1, Tier-2 capital) under Basel Norms, plus the Prompt Corrective Action (PCA) framework and due diligence --- so that a bank is not exposed to excessive risk.
\item Sustainability of business operations --- environmental-conservation norms require firms to mitigate environmental consequences.
\end{itemize}
\end{GovtPolicy}

\begin{GovtPolicy}
\begin{itemize}
\item Delhi air pollution / GRAP: when PM 2.5 and PM 10 exceed limits, construction activities are halted; not following such government regulations causes economic loss.
\item Level playing field: avoid anti-competitive practices (e.g.\ e-commerce deep discounts vs retail shops). The Competition Commission of India (CCI) protects against monopolies and maintains a level playing field.
\end{itemize}
\end{GovtPolicy}

\paragraph{Reducing Negative Externalities \& Inequality.}
\begin{KeyConcept}
Operations can cause environmental degradation and rising inequality, both of which government must contain because they ultimately hurt the market economy.
\end{KeyConcept}

\begin{ExampleBox}
\begin{itemize}
\item Great Economic Depression: if all companies take all profits and do not share with workers, no one can buy their products $\rightarrow$ recession $\rightarrow$ the market economy fails.
\item New Deal (USA): put money into people's hands through employment programmes and public expenditure to generate demand --- the ``poster boy of capitalism'' adopted socialist measures because the economy required it. Hence government must reduce inequality.
\item IndiGo crisis: to avoid market failure, government capped domestic fares (\rupee 34{,}000) when demand-supply broke down.
\end{itemize}
\end{ExampleBox}

\begin{QuickRevision}
Comply with regulations $\rightarrow$ more credibility, sustainability, risk management, level playing field, and reduced negative consequences \& inequality. ``Don't steal because of jail'' is weak; ``compliance protects the market economy that is in your own interest'' is strong.
\end{QuickRevision}

\subsection{Obligations Towards Society}
\begin{itemize}
\item A company is a product of society (it takes resources --- water, air, raw materials --- and leaves negative consequences on society). It is therefore always in debt to society and has obligations towards it.
\end{itemize}

\begin{KeyConcept}
\begin{itemize}
\item By contributing to society (e.g.\ healthcare and education), a company gains better employees and better consumers: good health/education $\rightarrow$ skilled human capital and consumers with better paying capacity $\rightarrow$ more demand and more productive workers.
\item Legitimacy / licence to operate: serving society earns the legitimacy to operate. Without it, a firm faces opposition and disruption.
\end{itemize}
\end{KeyConcept}

\begin{ExampleBox}
\begin{itemize}
\item Sterlite copper plant, Thoothukudi (Tamil Nadu): protests over environmental concerns $\rightarrow$ loss of legitimacy to operate.
\item Google data centres: even very large MNCs may not yet have earned local trust, and so may lack the legitimacy to operate, inviting opposition.
\end{itemize}
\end{ExampleBox}

\begin{PYQLink}
UPSC: How will CSR / corporate ethical governance make companies profitable? --- Argue through brand value, lower marketing expenditure, legitimacy to operate, employee commitment, productivity and motivation $\rightarrow$ profitability. When you are ethical you do no favour to anyone; you help only yourself.
\end{PYQLink}

\subsection{Ethical Concerns in Corporate Governance (India)}
\begin{KeyConcept}
\begin{itemize}
\item Importance of good corporate governance: raises investor confidence and foreign investment and lowers vulnerability to corruption.
\item Economic crises arise when morals fail and firms do not comply with corporate-governance norms --- Satyam scam (greed), Blue Smart Cab, the NPA crisis of banks, collusion (compliances skipped, nepotism shown).
\end{itemize}
\end{KeyConcept}

\paragraph{The Independent-Director Problem.}
\begin{CriticalPoint}
\begin{itemize}
\item The independent director exists to protect minority investors and provide a separate, in-house monitoring mechanism. But independent directors are often not truly independent.
\item Government PSUs: government holds a major stake and appoints ruling-party spokespersons as ``independent'' directors $\rightarrow$ favouritism, nepotism $\rightarrow$ they cannot monitor independently. ``Nepotism/favouritism in independent directors and boards can undermine the board's ability to provide unbiased oversight.''
\end{itemize}
\end{CriticalPoint}

\paragraph{Other Concerns.}
\begin{itemize}
\item Collusion (a huge risk --- e.g.\ bank NPAs), lack of oversight, lack of independence, nepotism/favouritism.
\end{itemize}

\begin{ExampleBox}
\begin{itemize}
\item Crony capitalism: a private company funds a political party's election campaign; after winning, the ruling party favours those companies in contracts and tenders --- unfair, and it destroys the level playing field.
\item Simply reverse the positive stakeholder points to generate concerns: glass ceiling, employee discrimination/exploitation, protest cases, fund misuse (Blue Smart), cheating consumers.
\item Further examples: Kingfisher Airlines (excessive risk), the NPA crisis, the Harshad Mehta scam, and ``front-hunt'' collusion.
\end{itemize}
\end{ExampleBox}

\subsection{Companies Act 2013 --- Provisions for Ethical Governance}

\paragraph{Women Director.}
\begin{GovtPolicy}
Mandatory to appoint at least one woman director in every listed company --- addressing the glass ceiling and ensuring representation.
\end{GovtPolicy}

\paragraph{Independent Directors.}
\begin{GovtPolicy}
\begin{itemize}
\item At least one-third of directors as independent directors --- purpose: protecting the interest of minority shareholders, improving corporate credibility, and fraud prevention \& detection. An in-house independent mechanism monitors the company and can detect fraud before it occurs. ``Independent'' means no dealings with the owners.
\item Misuse: appointing ruling-party spokespersons as independent directors of PSUs.
\end{itemize}
\end{GovtPolicy}

\paragraph{Audit Committee.}
\begin{GovtPolicy}
\begin{itemize}
\item Qualification: members must be able to read and understand financial statements --- a check so that not just anyone can be appointed.
\item Role: examination of financial statements and reports (verifying profit/loss/revenue are genuine, not faked); evaluation of internal financial control and risk management (the most important keyword); and monitoring the end-use of funds.
\item Analogy --- a bank minimises risk through KYC, CIBIL score, collateral and due diligence; the audit committee checks whether such risk-reduction systems are in place. It maintains corporate credibility and helps identify mistakes/scandals early.
\end{itemize}
\end{GovtPolicy}

\paragraph{Stakeholder Relationship Committee.}
\begin{GovtPolicy}
Purpose: to reduce and address grievances between shareholders and the board of directors.
\end{GovtPolicy}

\paragraph{Corporate Social Responsibility (CSR).}
\begin{DataFact}
\begin{itemize}
\item Eligible companies must spend 2\% of their (3-year average) net profit on society. India became the first country to legally mandate CSR.
\item Tata Steel spends 5--7\% of its profit after tax on CSR.
\end{itemize}
\end{DataFact}

\begin{ExampleBox}
\begin{itemize}
\item Tata Steel --- tribal areas: with many plants in tribal regions of Jharkhand and Odisha, Tata shows gratitude to tribal culture by establishing tribal cultural centres showcasing the legacy of tribes.
\item Atul initiative (Tata): trains persons with physical disability in the art of making recycled paper.
\item MANSI initiative --- Maternal and Newborn Survival Initiative (Tata Steel): aims to reduce child and infant mortality in states like Jharkhand; the MANSI project was awarded by the government.
\item Other positives: Patagonia (1\% of sales to environment); ITC e-Choupal (helping farmers).
\end{itemize}
\end{ExampleBox}

\begin{KeyConcept}
Consequences of ethics for a private organisation: employee satisfaction, better public image, licence to operate, and profitability.
\end{KeyConcept}

\subsection{Issues / Limitations of CSR}
\begin{itemize}
\item Additional corporate tax --- CSR is seen as a burden, so it is often not done in letter and spirit.
\item Projects not aligned to people's needs --- e.g.\ a company builds toilets where residents actually wanted a library; projects should follow the community's need, not the company's preference.
\item Dearth of well-meaning NGOs --- most credible NGOs are concentrated in urban and peri-urban areas, not in rural/remote regions, so projects are hard to execute where they are most needed.
\item Lack of a robust CSR policy --- firms lack a plan to spend large sums sustainably; CSR funds are viewed as a one-time disbursement with no long-term policy.
\item Others: corruption, duplication of activity, poor implementation and a skewed pattern of expenditure.
\end{itemize}

\begin{QuickRevision}
Corporate governance = fairness to every stakeholder (shareholders, consumers, employees, government/regulators, society).\\
Across all of them the refrain is ``good ethics is good business'' --- ethics helps the firm itself in the long run.\\
Companies Act 2013 tools: woman director, one-third independent directors, audit committee, stakeholder relationship committee, CSR (2\%).\\
Concerns: fake independence of directors, collusion, crony capitalism, NPAs, discrimination. CSR limits: burden mindset, mismatch with need, NGO scarcity, weak policy.
\end{QuickRevision}

\section{Ethics in International Relations}

\subsection{Background --- ``There is No Place for Principles in Politics''}
\begin{KeyConcept}
A common belief holds that there is no place for principles in politics --- politics is all about power: how to acquire it and how to remain powerful.
\end{KeyConcept}

\begin{ExampleBox}
\begin{itemize}
\item Machiavelli: ``the central goal of the ruler is to maintain and increase power,'' and the ends justify the means.
\item Chanakya / Kautilya: the doctrine of s\={a}m, d\={a}m, da\d{n}\d{d}a, bheda.
\item Bismarck (Unification of Germany): ``the great questions of the day will not be settled by speeches and majority decisions, but by iron and blood'' --- i.e.\ conflicts are settled by the policy of war, not democratic decision-making.
\item Donald Trump (to Zelensky): ``you don't have the cards'' --- implying a weak country is not in a position to argue; poor/weak countries never hold power.
\item Domestically the same is visible --- rising criminalisation, the role of money power, misuse of the Governor's office, and misuse of the anti-defection law to manufacture governments --- showing politics and principles not going hand in hand.
\end{itemize}
\end{ExampleBox}

\subsection{Why International Relations Are Different --- The Analogy}
A ladder of examples explains why regulation is weakest at the international level.

\begin{KeyConcept}
\begin{itemize}
\item Individual: if a person says ``I will do whatever I want,'' he cannot --- police, judiciary, law and the rule of law stand above him.
\item State government: if a state says ``I will be independent and follow nothing,'' it cannot --- the central government, Constitution, army and judiciary stand above it.
\item Country: if a country says ``I will do whatever I want,'' at the international level there is no corresponding enforcement power / rule of law. UN mechanisms are dysfunctional; the USA does much as it pleases, and no one stops it.
\end{itemize}
\end{KeyConcept}

\begin{CriticalPoint}
The crux: at the international level there is nothing as enforceable as domestic law. Enforcement is what matters, and it is missing --- hence the deep challenge in regulating countries' behaviour internationally.
\end{CriticalPoint}

\subsection{Why Countries Should Be Ethical --- The Arguments}
Since there is no international enforcer, the argument must be that it is in a country's own interest to be ethical. Each argument below pairs a reason with an example.

\paragraph{Global Nature of Problems.}
\begin{KeyConcept}
Terrorism, climate change, global warming, cyber-security and global pandemics (COVID) are global problems, and global problems require global solutions. No nation can tackle them alone $\rightarrow$ all share responsibility, which is the essence of ethics (duty beyond narrow self-interest). Violation example: the USA withdrawing from the Paris Climate Deal.
\end{KeyConcept}

\paragraph{To Resolve Conflicts Amicably \& Peacefully.}
\begin{KeyConcept}
Given the great inequality in military capabilities, if all conflicts are settled by military means the bigger fish always eats the smaller fish --- that suppresses conflict rather than solving it, so there is never lasting peace. ``Peace can never be kept by force''; ``the absence of tension does not mean peace.'' We must believe in peaceful resolution and a rule-based architecture.
\end{KeyConcept}

\paragraph{For Inclusive \& Sustainable Growth.}
\begin{KeyConcept}
Because of developmental inequalities, not all countries can fight all problems equally. India cannot say ``if Pakistan has polio, it is not our problem.'' Helping other countries helps ourselves --- e.g.\ providing vaccines/medicines to Africa, since it is better to contain a disease in its source region than to face it at home. The EU helps African nations tackle conflict and refugee flows in the source region.
\end{KeyConcept}

\paragraph{Overarching Rule-Based Architecture.}
\begin{KeyConcept}
National laws may not be sufficient --- a global law on terrorism is needed. Hence international conventions. (Afghanistan/Pakistan illustrate the terrorism gap.)
\end{KeyConcept}

\paragraph{To Have Flourishing Trade.}
\begin{KeyConcept}
Flourishing trade needs a level playing field and fair mechanisms to negotiate trade deals --- trade must not be enforced through bullying, or it will not last. Hence multilateral frameworks like the WTO (made dysfunctional by the USA, e.g.\ the paralysed dispute-settlement appellate mechanism).
\end{KeyConcept}

\begin{ExampleBox}
India \& TRIPS: India changed its IPR laws to move from process patents (a weaker form --- only the specific process is protected, so the product can be reverse-engineered by another process) to product patents (a stronger form --- the product itself cannot be made by any process). A concession made through international negotiation for the sake of flourishing trade.
\end{ExampleBox}

\paragraph{To Uphold Rule of Law \& Accountability.} Countries committing genocides against their own citizens must be held accountable --- upholding the international rule of law.

\paragraph{To Uphold Cultural Relativism \& Tolerance.}
\begin{KeyConcept}
Powerful nations tend to impose ``our culture is the best.'' (Trump to Zelensky: ``where is your suit?'' --- as if wearing a suit were a precondition for peace talks.) Without institutions like UNESCO promoting cultural diversity, the result is cultural imperialism --- powerful nations imposing their cultures. So we must promote diversity of culture and tolerance.
\end{KeyConcept}

\paragraph{Interdependence in the Age of Technology \& Globalisation.}
\begin{KeyConcept}
``In this age of technology and globalisation we are all interdependent on each other; the sooner we understand this, the better it is for all of us.'' A war in one place (e.g.\ Iran) affects everyone through energy supply, trade routes and global value chains. We must believe in humanity and collective well-being.
\end{KeyConcept}

\begin{CriticalPoint}
Common tone of all arguments: ethics benefits us; ethics is logical and rational. In effect you are telling every nation --- do not be selfish, or we all suffer.
\end{CriticalPoint}

\subsection{Which Ethical Principles to Follow in IR}
\begin{PYQLink}
Cosmopolitan Moral Theory (Immanuel Kant) --- a PYQ: every person is of equal moral worth; one cannot say American citizens are worth more than Africans. Therefore political/territorial boundaries are morally irrelevant --- where there is genocide or killing, borders do not excuse inaction. This guides us to uphold human rights and preserve humanity.
\end{PYQLink}
\begin{itemize}
\item Panchsheel principles --- respect for sovereignty and related tenets.
\item Solidarity and global cooperation --- the responsibility to help others.
\item Mutual respect for cultural diversity.
\item UN principles and human rights.
\end{itemize}

\subsection{Ethical Concerns in IR (Current-Affairs Examples)}
Human-rights violations, war crimes, environmental degradation, refugee crises, cyber-security problems and cultural imperialism --- in each you can show how international ethical principles are being violated.

\begin{QuickRevision}
IR lacks an enforcer (unlike domestic law) $\rightarrow$ argue ethics as self-interest.\\
Eight arguments: global problems, peaceful conflict resolution, inclusive growth, rule-based architecture, flourishing trade, rule of law/accountability, cultural relativism, interdependence.\\
Principles: Cosmopolitan Moral Theory (Kant), Panchsheel, solidarity \& cooperation, cultural respect, UN \& human rights.
\end{QuickRevision}

\subsection{Ethical Issues in International Funding}
Funding may be bilateral (country A to country B) or multilateral (IMF, World Bank, Asian Development Bank). On the surface it helps and promotes democracy --- but there can be hidden problems.

\begin{ExampleBox}
Hidden interests --- Kudankulam Nuclear Power Plant: the IB report stated the protests had hidden foreign-funded interests aimed at stalling India's development story / delaying the project --- funding arriving under the guise of environment or democracy but with a sinister design.
\end{ExampleBox}

\begin{CriticalPoint}
\begin{itemize}
\item Conditionality: country A funds country B on condition that, say, 70\% of the fund is spent buying equipment from country A --- reducing the recipient's sovereignty, even though it may have a cheaper, safer or different alternative.
\item Neo-imperialism: IMF/World Bank aid conditionalities (limit fiscal deficit, cut subsidies, reduce support programmes) promote only one, Western mode of development (liberalisation, private-sector-led growth, lower fiscal deficit) --- as if the recipient were not wise enough to have its own model. This restricts sovereignty (e.g.\ a party that promised free healthcare/education finds its options cut).
\end{itemize}
\end{CriticalPoint}

\begin{ExampleBox}
\begin{itemize}
\item Diversion of funds / terror funding: the Falah-e-Insaniat Foundation (FIF), operated by Hafiz Saeed (26/11), is cited as an NGO diverting funds to terror activities; corruption is a further concern.
\item Clinical trials: developed-country firms shift trials to poor countries to exploit cheaper labour, weak regulatory capacity and lower penalty risk. The Bhopal Gas Tragedy (1984) shows how consequences that would be enormously costly in the USA were not so in India.
\item China's Belt and Road Initiative (One Belt One Road): an example of hidden interest / debt-trap diplomacy.
\end{itemize}
\end{ExampleBox}

\begin{QuickRevision}
International-funding concerns: hidden interests (Kudankulam/IB report), conditionality (loss of sovereignty), neo-imperialism (IMF/WB, one Western model), diversion / terror funding (FIF--Hafiz Saeed), clinical trials (Bhopal), and debt-trap diplomacy (BRI).
\end{QuickRevision}

\section{Just War Theory}

\subsection{The Logic of Foreseeable Consequences}
\begin{KeyConcept}
\begin{itemize}
\item If you know the consequences of an action beforehand, you are held more responsible for it (e.g.\ smoking despite knowing it causes cancer; eating sugar despite diabetes).
\item Yet we still undertake certain actions despite negative consequences because they also carry positive consequences --- e.g.\ abortion (a life taken, but a mother's life saved) or law-and-order management (a lathi-charge or firing may cost lives, but saves lives and preserves order).
\end{itemize}
\end{KeyConcept}

\begin{CriticalPoint}
War is similar: it means people will be killed --- own soldiers, the enemy, and innocents. Because this is foreseeable, the decision-maker is more responsible and the decision must be taken with great care. War must be a method of last resort.
\end{CriticalPoint}

\subsection{Principles Before War (Jus ad Bellum)}
\begin{itemize}
\item \textbf{Last resort} --- war only after all peaceful alternatives are exhausted. In the Ramayana, Angad and Hanuman were sent first; in the Mahabharata, Krishna sought just five villages --- war came only when refused.
\item \textbf{Legitimate authority} --- war may be waged only by a legitimate authority (the state/government), which alone has the legitimacy and the duty to protect. Individuals waging war would be terrorists; Congress or BJP as parties cannot wage war.
\item \textbf{Reasonable assurance / probability of success} --- if the objective cannot be achieved, killing is unjustified. Terrorists have no probability of success, which is one reason their actions are unethical.
\item \textbf{Right intention} --- the objective of a just war is to re-establish peace; the intention must be just (justice, not image-protection or chest-thumping). Kant stressed intention above all.
\item \textbf{Proportionality} --- resorting to war must lead to more good than harm; violence should escalate gradually (lathi-charge $\rightarrow$ water cannon $\rightarrow$ tear gas), not straight to tanks. India's use of force is gradual, non-escalatory and voluntarily self-restrained.
\end{itemize}

\begin{ExampleBox}
Legitimacy to decision-making --- all-party meeting: convened to lend legitimacy to a decision on the use of force, keeping the response non-escalatory.
\end{ExampleBox}

\subsection{During War (Jus in Bello) \& After War (Jus post Bellum)}
\begin{KeyConcept}
\begin{itemize}
\item During war: minimise civilian harm and risk; commit no war crimes (e.g.\ mass rape, killing of civilians).
\item Post-war: ensure justice, rehabilitation and accountability mechanisms.
\end{itemize}
\end{KeyConcept}

\subsection{Purpose \& Exam Application}
\begin{PYQLink}
UPSC --- ``War is also a tool of diplomacy'': not \emph{every} war qualifies; only a just war that follows these ethical principles can be a tool of diplomacy.
\end{PYQLink}

\begin{CriticalPoint}
Just War Theory does not advocate or glorify war. It exists to explain when war can be ethical, because the decision to kill is grave and must be reasoned. War, like abortion, is inherently a grave wrong undertaken only under strict justification.
\end{CriticalPoint}

\begin{ExampleBox}
1971 War: had it not happened, the refugee crisis would have been far worse --- a case where war, as a last resort, produced a just result. The Ramayana war is similarly cited.
\end{ExampleBox}

\begin{QuickRevision}
Foreseeable killing makes the war decision grave $\rightarrow$ last resort.\\
Before war (jus ad bellum): last resort, legitimate authority, probability of success, right intention (re-establish peace), proportionality (more good than harm).\\
During war: minimise civilian harm, no war crimes. After war: justice, rehabilitation, accountability.\\
The theory explains ethical war; it does not promote war. Just war = a tool of diplomacy (1971).
\end{QuickRevision}
